\chapter{放射防护措施评价}

\section{场所布局与分区}

\subsection{布局合规核查}

\begin{table}[H]
\centering
\begin{tabular}{|c|c|c|c|}
\hline
项目 & 设计情况 & 国标 GBZ121-2020 要求 & 评价 \\
\hline
机房位置 & 肿瘤楼一层端部 & 放疗机房宜建筑底层端部 & 符合 \\
\hline
分区设置 & 机房 + 迷道 = 控制区；周边走道 / 准备间 = 监督区 & 分控制、监督两区 & 符合 \\
\hline
控制室 & 机房西侧独立设置 & 控制室与机房物理分隔 & 符合 \\
\hline
迷道配置 & 直型迷路入口 & γ 后装机房必须设置迷道 & 符合 \\
\hline
\end{tabular}
\end{table}


\subsection{分区划分}

\begin{itemize}
\item \textbf{控制区}：机房室内 + 迷道，门禁管控、警示标识，非工作人员禁止入内；
\item \textbf{监督区}：机房外侧过道、准备间，定期辐射巡检。
\end{itemize}

\subsection{机房尺寸}

机房净尺寸：长 5.05m× 宽 4.80m× 高 4.00m，使用面积 24.24㎡，空间满足设备安装与临床操作，符合规范。

\section{机房屏蔽设计与核算}


\begin{table}[H]
\centering
\caption{机房屏蔽构造（混凝土密度≥2.35t/m³，全部利旧）}
\begin{tabular}{|c|c|c|c|}
\hline
屏蔽部位 & 墙体厚度 & 材质 & 备注 \\
\hline
东墙 / 南墙 / 北墙 / 顶板 & 800mm & 钢筋混凝土 & 利旧 \\
\hline
西墙（靠控制室） & 1700mm & 钢筋混凝土 & 利旧 \\
\hline
防护大门 & 10mmPb 铅防护电动门 & 铅钢复合 & 利旧 \\
\hline
地面 & 原生夯实土层 & 素土 & 无额外屏蔽 \\
\hline
\end{tabular}
\end{table}



\begin{table}[H]
\centering
\caption{各关注点剂量控制目标}
\begin{tabular}{|c|c|c|c|c|c|c|c|c|c|}
\hline
关注点 & 位置 & 人员类型 & 周控制剂量 & 居留因子 T & 使用因子 U & 周工时 t & 核算$H_c$ & 限值$H_{c,max}$ & 最终控制$H_c$ \\
\hline
a & 准备间 & 工作人员 & 5μSv & 1/2 & 1 & 5h & 2.00μSv/h & 10 & 2.00 \\
\hline
b & 狭窄过道 & 公众 & 5μSv & 1/40 & 1 & 5h & 40.00μSv/h & 10 & 10.00 \\
\hline
c & 预留加速器区 & 公众 & 5μSv & 1/8 & 1 & 5h & 8.00μSv/h & 10 & 8.00 \\
\hline
d & 控制室 & 工作人员 & 100μSv & 1 & 1 & 5h & 20.00μSv/h & 2.5 & 2.50 \\
\hline
e & 屋面平台 & 公众 & 5μSv & 1/40 & 1 & 5h & 40.00μSv/h & 10 & 10.00 \\
\hline
f & 防护门外 & 公众 & 5μSv & 1/8 & 1 & 5h & 8.00μSv/h & 10 & 8.00 \\
\hline
\end{tabular}
\end{table}



\begin{table}[H]
\centering
\caption{屏蔽核算结果汇总}
\begin{tabular}{|c|c|c|c|c|c|}
\hline
点位 & 所需理论厚度 & 实际设计厚度 & 墙外计算剂量 (μSv/h) & 管控限值 & 评价 \\
\hline
a 准备间 & 420mm & 800mm & 6.33×10⁻³ & 2 & 符合 \\
\hline
b 过道 & 434mm & 800mm & 3.89×10⁻² & 10 & 符合 \\
\hline
c 南侧预留机房 & 419mm & 800mm & 2.49×10⁻² & 8 & 符合 \\
\hline
d 控制室 & 543mm & 1700mm & 5.08×10⁻² & 2.5 & 符合 \\
\hline
e 屋顶 & 335mm & 800mm & 8.61×10⁻³ & 10 & 符合 \\
\hline
f 防护门外 & 5mmPb & 10mmPb & 0.313 & 10 & 符合 \\
\hline
\end{tabular}
\end{table}


\begin{quote}
结论：所有屏蔽厚度＞理论计算需求，墙外辐射剂量全部低于控制限值，屏蔽满足国标。
\end{quote}

\section{安全防护设施}


\begin{table}[H]
\centering
\caption{安全装置配置核查}
\begin{tabular}{|c|c|c|c|}
\hline
设施名称 & 配置情况 & 标准要求 & 结论 \\
\hline
固定式室内辐射报警器 & 探测器机房内，主机控制室 & 后装机房必须安装固定式剂量报警 & 符合 \\
\hline
门机联锁 & 关门才可送源，开门自动回源停机 & 防护门未关闭禁止出源 & 符合 \\
\hline
机房内外开门按钮 & 室内、门外双应急开门 & 机房内设手动开门装置、防挤压 & 符合 \\
\hline
电离辐射标识 + 工作指示灯 & 防护门张贴警示，亮灯禁止入内 & 出入口设置警告标识与工作指示灯 & 符合 \\
\hline
多点位急停按钮 & 设备本体、机房多面墙体、控制台均设急停 & 全室可见、随手可触紧急停机 & 符合 \\
\hline
视频 + 对讲系统 & 机房全视角监控、医患双向对讲 & 控制室实时观察患者、双向通话 & 符合 \\
\hline
应急储源罐 + 长柄镊子 & 库房备用，卡源故障应急存源 & γ 后装项目必备应急储源设施 & 符合 \\
\hline
手动回源工具 & 配套备用，自动故障手动回源 & 自动回源失效可手动回收放射源 & 符合 \\
\hline
\end{tabular}
\end{table}


\subsection{个人防护用品}

配备 2 台个人辐射报警仪；放疗禁止患者家属陪检。

\subsection{三废}

无放射性固 / 气 / 废水；废旧$^{192}\text{Ir}$废源由原厂上门回收处置。

\subsection{机房通风}

机房容积 140m³，机械排风设计换气 7 次 /h（国标≥4 次 /h）；上进风、下部侧墙排风，排风管道直通屋顶高空排放，用于消除臭氧、氮氧化物。

\subsection{管线穿墙}

电缆沟、通风管道采用多折弯管穿墙，规避高居留控制室，防止射线沿孔洞泄漏。

\section{小结}

机房分区合理、屏蔽厚度富余、安全联锁 / 报警 / 通风 / 应急器材配置齐全，全部满足 GBZ121-2020。
